\section{Classical sequent calculus}
\label{sec:prop-deduction-02}

\begin{definition}[Sequent]
A \emph{finite multiset of formulas} is a function
\(\Gamma:\Fm_{\CPL}\to\N\) whose support
\[
 \operatorname{supp}(\Gamma)
 \defeq\{\varphi\in\Fm_{\CPL}:\Gamma(\varphi)>0\}
\]
is finite.  The empty multiset is the constant-zero function.  If
\(\Gamma\) and \(\Pi\) are finite multisets, then the comma
\(\Gamma,\Pi\) denotes their pointwise sum; a displayed formula \(\varphi\)
in a comma list denotes the multiset having value \(1\) at \(\varphi\) and
\(0\) elsewhere.  Thus \(\Gamma,\varphi,\varphi\) increases the multiplicity
of \(\varphi\) by two.  In the first-order sequent chapter the same definition
is used with \(\Fm(\Lang)\) in place of \(\Fm_{\CPL}\).

A \emph{sequent} is an expression \(\Gamma\Rightarrow\Delta\), where
\(\Gamma\) and \(\Delta\) are finite multisets of formulas.  It is satisfied
by a valuation iff either some formula in \(\operatorname{supp}(\Gamma)\) is
false or some formula in \(\operatorname{supp}(\Delta)\) is true.  With
respect to satisfaction, an empty antecedent therefore behaves as
\(\topv\) and an empty succedent as \(\bottom\).

When a finite multiset occurs to the left of \(\entails\) or \(\derives\),
this document defines
\[
 \Gamma\entails\varphi
 \quad\Longleftrightarrow\quad
 \operatorname{supp}(\Gamma)\entails\varphi,
 \qquad
 \Gamma\derives_{\mathcal C}\varphi
 \quad\Longleftrightarrow\quad
 \operatorname{supp}(\Gamma)\derives_{\mathcal C}\varphi.
\]
Multiplicity is therefore proof-structural data in a sequent, while semantic
consequence and natural-deduction assumptions depend only on the support.
\end{definition}

The calculus \(\LK_{\CPL}\) has two forms of \emph{initial sequent}:
identity sequents \(\Gamma,\varphi\Rightarrow\varphi,\Delta\) and
\(\bottom\)-sequents \(\Gamma,\bottom\Rightarrow\Delta\).  It also has
the following rules.  Exchange is
implicit because contexts are multisets.
\[
\begin{array}{c@{\qquad}c}
\Rule[\(W_L\)]{\Gamma\Rightarrow\Delta}{\Gamma,\varphi\Rightarrow\Delta}
&
\Rule[\(W_R\)]{\Gamma\Rightarrow\Delta}{\Gamma\Rightarrow\Delta,\varphi}
\\[2.4ex]
\Rule[\(C_L\)]{\Gamma,\varphi,\varphi\Rightarrow\Delta}
 {\Gamma,\varphi\Rightarrow\Delta}
&
\Rule[\(C_R\)]{\Gamma\Rightarrow\Delta,\varphi,\varphi}
 {\Gamma\Rightarrow\Delta,\varphi}
\end{array}
\]
\[
\begin{array}{c@{\qquad}c}
\Rule[\(\neg L\)]{\Gamma\Rightarrow\Delta,\varphi}
 {\Gamma,\neg\varphi\Rightarrow\Delta}
&
\Rule[\(\neg R\)]{\Gamma,\varphi\Rightarrow\Delta}
 {\Gamma\Rightarrow\Delta,\neg\varphi}
\\[2.6ex]
\Rule[\(\land L\)]{\Gamma,\varphi,\psi\Rightarrow\Delta}
 {\Gamma,\varphi\land\psi\Rightarrow\Delta}
&
\Rule[\(\land R\)]{
 \Gamma\Rightarrow\Delta,\varphi\qquad
 \Gamma\Rightarrow\Delta,\psi}
 {\Gamma\Rightarrow\Delta,\varphi\land\psi}
\\[2.6ex]
\Rule[\(\lor L\)]{
 \Gamma,\varphi\Rightarrow\Delta\qquad
 \Gamma,\psi\Rightarrow\Delta}
 {\Gamma,\varphi\lor\psi\Rightarrow\Delta}
&
\Rule[\(\lor R\)]{\Gamma\Rightarrow\Delta,\varphi,\psi}
 {\Gamma\Rightarrow\Delta,\varphi\lor\psi}
\\[2.6ex]
\Rule[\(\to L\)]{
 \Gamma\Rightarrow\Delta,\varphi\qquad
 \Gamma,\psi\Rightarrow\Delta}
 {\Gamma,\varphi\to\psi\Rightarrow\Delta}
&
\Rule[\(\to R\)]{\Gamma,\varphi\Rightarrow\Delta,\psi}
 {\Gamma\Rightarrow\Delta,\varphi\to\psi}.
\end{array}
\]
The cut rule is
\[
 \Rule[Cut]{\Gamma\Rightarrow\Delta,\varphi\qquad
             \Sigma,\varphi\Rightarrow\Pi}
            {\Gamma,\Sigma\Rightarrow\Delta,\Pi}.
\]
A \(\LK_{\CPL}\)-derivation is a finite tree whose leaves are initial
sequents and whose internal nodes are licensed rule instances.

\begin{example}[A sequent derivation]
The sequent \(\Rightarrow P\to(P\lor Q)\) has the cut-free derivation
\[
 \Rule[\(\to R\)]{
   \Rule[\(\lor R\)]{P\Rightarrow P,Q}{P\Rightarrow P\lor Q}}
  {\Rightarrow P\to(P\lor Q)}.
\]
The upper sequent is an identity sequent.
\end{example}

